% ch17-reflection.tex — פרק אישי: כל אחד מהמגישים כותב את הסעיף שלו, בקול שלו.
% הפרק לא נכנס לדוח עד ששני הסעיפים מלאים. כשהם מלאים: להוסיף ב-main.tex אחרי ch16:
%   % ch17-reflection.tex — פרק אישי: כל אחד מהמגישים כותב את הסעיף שלו, בקול שלו.
% הפרק לא נכנס לדוח עד ששני הסעיפים מלאים. כשהם מלאים: להוסיף ב-main.tex אחרי ch16:
%   % ch17-reflection.tex — פרק אישי: כל אחד מהמגישים כותב את הסעיף שלו, בקול שלו.
% הפרק לא נכנס לדוח עד ששני הסעיפים מלאים. כשהם מלאים: להוסיף ב-main.tex אחרי ch16:
%   % ch17-reflection.tex — פרק אישי: כל אחד מהמגישים כותב את הסעיף שלו, בקול שלו.
% הפרק לא נכנס לדוח עד ששני הסעיפים מלאים. כשהם מלאים: להוסיף ב-main.tex אחרי ch16:
%   \input{chapters/ch17-reflection}
% כל אחד עושה commit ו-push מהחשבון שלו, עם הודעה כמו:
%   "Add Paz's reflection: why this topic and what I learned"
%
% שאלות מנחות (לא חובה לענות על כולן, 150 עד 250 מילים לכל אחד):
%   1. למה דווקא הנושא הזה? מה ראית בעבודה, בפיד או בחדשות שגרם לך לבחור בו?
%   2. מה הפתיע אותך במחקר? (למשל: שכבר יש מתחרים, ש-C2PA לא שורד צילום מסך)
%   3. איפה הסוכן טעה ואיך תפסת את זה?
%   4. מה היית עושה אחרת אם היית מתחיל היום?
%   5. מה מזה תיקח לעבודה או ללימודים?
% מותר לכתוב בגוף ראשון, בשפה פשוטה. זה בדיוק הקול האישי שהמרצה ביקש.

\chapter{מה למדנו: בקול שלנו}\label{ch:reflection}

\section{גיא יערי}
% גיא כותב כאן.

\section{פז ניסים}
% פז כותב כאן.

% כל אחד עושה commit ו-push מהחשבון שלו, עם הודעה כמו:
%   "Add Paz's reflection: why this topic and what I learned"
%
% שאלות מנחות (לא חובה לענות על כולן, 150 עד 250 מילים לכל אחד):
%   1. למה דווקא הנושא הזה? מה ראית בעבודה, בפיד או בחדשות שגרם לך לבחור בו?
%   2. מה הפתיע אותך במחקר? (למשל: שכבר יש מתחרים, ש-C2PA לא שורד צילום מסך)
%   3. איפה הסוכן טעה ואיך תפסת את זה?
%   4. מה היית עושה אחרת אם היית מתחיל היום?
%   5. מה מזה תיקח לעבודה או ללימודים?
% מותר לכתוב בגוף ראשון, בשפה פשוטה. זה בדיוק הקול האישי שהמרצה ביקש.

\chapter{מה למדנו: בקול שלנו}\label{ch:reflection}

\section{גיא יערי}
% גיא כותב כאן.

\section{פז ניסים}
% פז כותב כאן.

% כל אחד עושה commit ו-push מהחשבון שלו, עם הודעה כמו:
%   "Add Paz's reflection: why this topic and what I learned"
%
% שאלות מנחות (לא חובה לענות על כולן, 150 עד 250 מילים לכל אחד):
%   1. למה דווקא הנושא הזה? מה ראית בעבודה, בפיד או בחדשות שגרם לך לבחור בו?
%   2. מה הפתיע אותך במחקר? (למשל: שכבר יש מתחרים, ש-C2PA לא שורד צילום מסך)
%   3. איפה הסוכן טעה ואיך תפסת את זה?
%   4. מה היית עושה אחרת אם היית מתחיל היום?
%   5. מה מזה תיקח לעבודה או ללימודים?
% מותר לכתוב בגוף ראשון, בשפה פשוטה. זה בדיוק הקול האישי שהמרצה ביקש.

\chapter{מה למדנו: בקול שלנו}\label{ch:reflection}

\section{גיא יערי}
% גיא כותב כאן.

\section{פז ניסים}
% פז כותב כאן.

% כל אחד עושה commit ו-push מהחשבון שלו, עם הודעה כמו:
%   "Add Paz's reflection: why this topic and what I learned"
%
% שאלות מנחות (לא חובה לענות על כולן, 150 עד 250 מילים לכל אחד):
%   1. למה דווקא הנושא הזה? מה ראית בעבודה, בפיד או בחדשות שגרם לך לבחור בו?
%   2. מה הפתיע אותך במחקר? (למשל: שכבר יש מתחרים, ש-C2PA לא שורד צילום מסך)
%   3. איפה הסוכן טעה ואיך תפסת את זה?
%   4. מה היית עושה אחרת אם היית מתחיל היום?
%   5. מה מזה תיקח לעבודה או ללימודים?
% מותר לכתוב בגוף ראשון, בשפה פשוטה. זה בדיוק הקול האישי שהמרצה ביקש.

\chapter{מה למדנו: בקול שלנו}\label{ch:reflection}

\section{גיא יערי}
% גיא כותב כאן.

\section{פז ניסים}
% פז כותב כאן.
